\chapter{Foreward to the First Edition}

Yeah, that's right, this damn thing is big enough that it gets its own
foreward.  Go ahead and skip right over it... I \emph{dare} you.

Ahhem... Err... Undergrad sucks.  Going back through these problems I
can't help but notice that undergrad in the oughties doomed many
students to failure while simultaneously relegating others to an
unimaginative future.  In so many cases (check out
\outref{1}{24}{theQuestion}) the \emph{cool} or even the natural way
to work a problem is pretty different than what the author imagines.
It's not that they were doing a ``bad'' job, so much as they didn't
have a better approach available.  Sometimes the cool way requires
tools that they haven't taught, or it takes time that you don't have,
or ...  However, prior teaching methodologies lacked a tool that we
have available to us today: the \ac{llm}.  Additionally, in undergrad
I at least didn't do a particularly good job of focusing on skill
development and instead focused primarily upon expanding the number of
connections in my cognitive network.  As a precursor to exploring
quantum computation, I decided to refresh my basics in quantum
mechanics so I worked all the problems in chapter 1 of Sakurai.

So, if I used an \ac{llm}, did I really \emph{solve} the problems
myself?  That's honestly a question I'm truly worried about.
Unfortunately, I'm doing this almost entirely unsupported.  I have no
study group, \ac{ta}, or even so much as a pen-pal to work through
these with.  Left to my own devices, my solutions would have been
\emph{vastly} less interesting.  There are integrals in the Ice Pick
problem (\outref{1}{24}{theQuestion}) that I would not have been able
to take on my own.  I went back through my prob stats book from
undergrad and even knowing what to look for I would \emph{not} have
made the connection to $\chi^2$.\footnote{In fairness, I have recently
come to the conclusion that I think the Hogg and Tanis book sucks and
have since ordered Jaynes.}  Is it OK to use an integral table?  What
about Wolfram Alpha?  What about a calculator?  It's going to be
years, if not decades, before we know the consequences on human
cognition and learning and creativity of heavy use of \acp{llm}.  To
guard against these hypothesized (and in some cases observed)
failure-modes, I worked quite hard to solve everything on my own and
use the \ac{llm} largely to check my mathematical rigor and every now
and then to provide a hint.  In exactly one case, specifically the
``standard'' (read: ``lame'') way of solving
\outref{1}{34}{theAnswer}, I used the method that the \ac{llm}
recommended.  That was included because I figure I should include the
``standard'' solution in all of the questions, which I do in all
problems though I usually just stumble across it while doing it
the... umm... cool way.  You also \emph{constantly} have to tell the
damn things not to solve the problem for you... it's like having a
kid.

I tried to keep the material fun to read, and to provide enough level
of detail that the slow kids (such as yours truly) can keep up.  I
uh... I also may have gone a little overboard on the Ice Pick
problem...  You'll see...  I show you many of the dead-ends that I
explored and traps where it's easy to make a mistake and I try not to
be shy about calling myself a dumbass... don't worry... I'm sure
you'll agree with me before too long.

The most useful advice I've ever received about grad school is to view
it as trade school -- you're learning the trade of being a physicist.
It's not reasonable for us to have to re-derive everything, so I try
pretty hard to help the reader stand on the shoulders of the
proverbial giants.  There are a \emph{lot} of tools and I spent a lot
of time just practicing with the tools... and sometimes breaking
one... I promise I'll buy you a new one.  This kind of trade-shcool
approach was incredibly effective, especially when coupled with an
\ac{llm}.  After all, you don't learn physics by banging your head
against an integral that everyone else has already solved because they
know ``the trick''.

Jocosity aside, I sincerely hope that this guide helps you in your
journeys... unless you're a Krasnov supporter... in which case I have
a bridge to sell you and would recommend you drop out of your physics
curriculum and get your real estate license.

Sarcastically,

Harrison
